\section{Abstract}

Agricultural subsurface (tile) drainage underpins crop production on poorly drained and irrigated land, yet the same conduits that make these soils farmable also short-circuit the landscape's natural buffering, delivering nitrate, phosphorus, salts, and a widening list of constituents to ditches, streams, and coastal waters. Decades of field experiments, monitoring, and modelling have produced a detailed literature on drained-system behaviour and on the practices---nutrient management, controlled drainage, cover crops, edge-of-field treatment, salinity reclamation---that might reconcile productivity with water quality \parencite{Randall2001}. This review is the research-frontier component of a larger synthesis, framed deliberately around what remains uncertain, untested, or unstudied rather than around settled results. Its recurring meta-finding: ideas have outrun the data able to test them--- headline results rest disproportionately on short records, single sites, and narrow agro-ecological windows, and progress in every domain is bounded by the length, density, replication, and consistency of the underlying observations. We proceed from enabling to domain frontiers, then to consequences and integration. The enabling foundations are an observational base of longer, denser, and harmonised monitoring; process-based drainage and water-quality models whose credibility is limited less by their formulation than by the narrow envelope in which they have been validated; and the coupled problems of mapping the largely unmapped drainage network and scaling instrumented-field knowledge to the watersheds where water-quality targets are set. On this footing we treat the domain frontiers of nitrogen, phosphorus, and cropping-system losses; active water-table management and drainage-system redesign; edge-of-field treatment such as bioreactors, wetlands, and phosphorus-removal structures; and salinity reclamation, drainage-water reuse, and disposal. Because removing a constituent from one pathway can shift it to another and altering the field water regime can trigger slow biogeochemical and pedological change, we then consolidate the unintended consequences and multi-objective tradeoffs---greenhouse-gas coupling, pollution swapping, emerging contaminants, trace-element legacies, and drainage-induced soil evolution---that nitrate- and yield-centred accounting has overlooked. Finally, we draw the integrative strands together: no single practice or model is a silver bullet, and decisive advances will come from coupling mechanistic understanding with spatial targeting, economics, policy, and shared decision-support infrastructure. The unifying message is that the next decade of progress will be secured less by new equations or novel practices than by sustained, well-instrumented, openly reported, and geographically diverse field campaigns, coordinated data standards, and interdisciplinary integration tying mechanism to tools and tools to adoption.
